\documentclass{report}
\usepackage{amsmath,amssymb}
\setlength{\parindent}{0mm}
\setlength{\parskip}{1em}
\begin{document}
\begin{center}
\rule{6in}{1pt} \
{\large Gary W. Howell \\
{\bf LU Preconditioning for Iterative Solution of Overdetermed Sparse Least Squares Problems }}

512 Farmington Woods Dr Cary NC 27511
\\
{\tt gwhowell@ncsu.edu}\\
Marc Baboulin \end{center}

We investigate how to use an $LU$ factorization with the classical {\tt
lsqr} routine for solving overdetermined sparse least squares problems.
Usually $L$ is much better conditioned than $A$. Thus iterating with $L$
instead of $A$ results in faster convergence.
Numerical experiments using Matlab illustrate the good behavior of our
algorithm in terms of storage and convergence.
This paper explores a preliminary shared memory implementation using
SuperLU factorization.


We consider the overdetermined full rank LLS problem
\begin{equation}
\label{eq:llsp}
\min_{x \in \mathbb{R}^n} \neucs{Ax - b},
\end{equation}
with $A \in \mathbb{R}^{m \times n}, m \geq n$ and $b \in \mathbb{R}^m$.
When $A$ is sparse, direct methods based on $QR$ factorization or the
normal equations are not always suitable because the $R$ factor or $A^TA$
can be dense.
A common iterative method to find the least squares solution $x$
is to solve the normal equations

\begin{equation}
\label{eq:ne}
A^TAx = A^Tb,
\end{equation}
by applying the conjugate gradient (CG) algorithm to $A^TA$. In this case the matrix
$A^T A$ does not need to be explicitly formed, avoiding possible fill-in
in the formation of $A^TA$.

with other sparse linear systems, preconditoning techniques based on
incomplete factorizations can improve convergence. One method to
precondition the normal equations~(\ref{eq:ne}) is to perform an
incomplete Cholesky decomposition of $A^T A$ (e.g., RIF preconditioner,
Benzi, 2003).

When $A^TA$ and its Cholesky factorization are denser than $A$, it is
natural to wonder if the $LU$ factorization of $A$ can be used in solving
the least squares problem. In this paper we use an $LU$ factorization of
the rectangular matrix $A=\left(
\begin{array}{c}
A_1\\
A_2\\
\end{array}
\right)
$ where $L$ is unit lower trapezoidal and $U$ is upper triangular.
For the nonpivoting case, the normal equations~(\ref{eq:ne}) become
\begin{equation}
\label{eq:neL}
L^TLy=c,
\end{equation}
with $c = L^T b, U x = y $, and we can apply CG iterations on~(\ref{eq:neL}).

Least squares solution using $LU$ factorization has been explored by
several authors. Peters and Wilkinson (1970) and Bj\"orck and Duff (1980)
give direct methods. This work follows Bj\"orck and Yuan (1999) using
conjugate gradient methods based on $LU$ factorization, an approach worth
revisiting because of the recent progress in sparse $LU$ factorization.
The lsqrLU (2015) algorithm presented here uses a lower trapezoidal $L$
returned from a direct solver package. Here we use an $LU = PAQ$
factorization from shared memory SuperLU (X. LI 2005), comparing
iteration with $L$ to the $U^{-1}A$ iteration suggested by Saunders.
Because other direct solver packages offer scalable sparse $LU$
factorizations, it appears likely that the algorithm used here can also
be used to solve larger problems.
The rate of linear convergence for CG iterations on the normal equations is
$$K = \frac {\kappa - 1 }{\kappa + 1},$$
where $\kappa=\sqrt{\cond(A^TA)}=\cond(A)$ and $\cond(A)$ denotes the
2-norm condition number of $A$ (ratio of largest and smallest singular
values of $A$).
In our experiments, $L$ is often much better conditioned than $A$, so
convergence of the CG method is relatively rapid. Moreover, the total
number of nonzeros
in $L$ and $U$ is usually less than in the sparse Cholesky factorization of
$A^T A$.


\end{document}
